\section{Related Work}
\label{sec:related}

\subsection{Meta-learners for CATE}
\label{sec:related:xlearner}
Meta-learners decompose CATE estimation into a sequence of regression
problems, each solved by an off-the-shelf supervised learner. The
S-Learner fits a single outcome model $\mu(x, w)$ and takes
$\tau(x) = \mu(x, 1) - \mu(x, 0)$; the T-Learner fits separate
$\mu_0, \mu_1$ and takes their difference. Both are consistent but
suffer from regularisation bias: the S-Learner under-fits $\tau$
when $w$ is a weak predictor of $y$; the T-Learner over-fits each
arm when one arm is much smaller than the other. The X-Learner of
\citet{kunzel2019metalearners} repairs this by imputing
counterfactual outcomes and regressing on the resulting pseudo-effects,
then weighting the two fits by propensity. The construction is
particularly well-suited to treated/control imbalance, which is the
norm outside of randomised experiments.

\citet{kennedy2020optimal} shows that replacing the X-Learner's
imputed pseudo-outcomes with \emph{doubly robust} (DR) pseudo-outcomes
$D_1 = \mu_1 - \mu_0 + (Y - \mu_1)/\pi$ yields a target whose
second-stage regression consistently estimates $\tx$ under a rate
condition on the nuisance fits that is weaker than either the T- or
X-Learner requires separately. The R-Learner of \citet{nie2021quasi}
takes a complementary quasi-oracle route via residualisation.
\citet{curth2021nonparametric} survey the meta-learner landscape and
formalise a theory of conditions under which each
variant is preferred; a follow-up \citep{curth2023search} argues that
much of the apparent model-selection variance in the CATE literature
is benchmark-dependent rather than fundamental. Our method uses the
DR pseudo-outcome target --- the imputation and debiasing happen in
Phase 2 of our pipeline (Section~\ref{sec:method}) --- and inherits
its rate properties.
What classical meta-learners leave open is uncertainty quantification
on $\tx$: the standard recipe is bootstrap, which inherits the
sensitivity of the base learner to outliers.

\subsection{Bayesian causal inference and MCMC}
\label{sec:related:bayesian}
Bayesian causal inference on $\tx$ has been dominated by flexible
tree-based priors. BART itself \citep{chipman2010bart} provides the
backbone; Causal BART \citep{hill2011bayesian} adapts it to the
potential-outcome framework, producing a posterior over the
counterfactual response surfaces from which $\tx$ is derived. BCF
\citep{hahn2020bcf} decouples the treatment effect from the
prognostic surface using separate BART priors and yields substantial
improvements on heterogeneous response; SBCF
\citep{caron2022shrinkage} extends BCF with horseshoe-style global-
local shrinkage on the treatment-effect terms for sparse $\tx$.
Alternative Bayesian nonparametric formulations include Gaussian
processes \citep{alaa2018gp} and neural-network architectures with
uncertainty estimation \citep{shi2019adapting, jesson2021uncertainty}.
A recent line \citep{zheng2022bayesian} extends Bayesian CATE to
multi-treatment and partial-identification settings.
Frequentist semiparametric alternatives --- generalised random
forests \citep{athey2019generalized}, double/debiased machine
learning \citep{chernozhukov2018double}, and conformal inference for
individual treatment effects
\citep{lei2021conformal, alaa2023conformal, jin2023sensitivity} ---
deliver calibrated intervals through different routes (asymptotic
normality, orthogonal scores, finite-sample distribution-free
guarantees) but produce point-estimate-plus-interval objects rather
than full posteriors over $\tx$.
Modern probabilistic-programming languages (NumPyro, Stan, PyMC)
make Bayesian architectures easy to implement, and have been used to
build custom CATE models with hierarchical priors or non-conjugate
likelihoods.

The common thread is a Gaussian outcome likelihood on the target
scale: either directly on $Y$ (Causal BART) or on the transformed
pseudo-outcome (most custom constructions). Under heavy tails this
assumption distorts the posterior: a single whale's squared residual
dominates the likelihood gradient, pulling the posterior mean toward
the contaminated point and inflating the posterior variance. Our
contribution is to replace the Gaussian layer with a Welsch
redescending pseudo-likelihood while keeping the rest of the
meta-learner + MCMC architecture standard.

\subsection{Robust statistics for heavy-tailed regression}
\label{sec:related:robust}
\citet{huber1964robust} introduced the minimax framework for
point estimation under contaminated location models, proving that
the Huber $\psi$-function $\psi_\delta(r) = r \cdot
\mathbf{1}(|r| \le \delta) + \delta \cdot \mathrm{sign}(r) \cdot
\mathbf{1}(|r| > \delta)$ is minimax-optimal over the
$\epsilon$-contamination neighbourhood $\mathcal{F}_\epsilon$. The
optimal tuning constant $\delta$ is the unique solution of
\begin{equation}
  \frac{\phi(\delta)}{\delta} - (1 - \Phi(\delta))
  \;=\; \frac{\epsilon}{2(1 - \epsilon)},
  \label{eq:huber_minimax}
\end{equation}
where $\phi, \Phi$ denote the standard normal PDF and CDF. Classical
values are $\delta = 1.345$ at $\epsilon = 5\%$ and $\delta = 0.5$
at $\epsilon \approx 40\%$; the full table and a numerical
derivation appear in Appendix~\ref{app:huber_derivation}. The
asymptotic relative efficiency (ARE) of Huber estimation vs
maximum-likelihood under a truly Gaussian model follows from
$E[\psi_\delta^2] / E[\psi'_\delta]^2$ and is $\approx 95\%$ at
$\delta = 1.345$ and $\approx 79\%$ at $\delta = 0.5$. Redescending
M-estimators push robustness further by letting the influence
function return smoothly to zero for very large residuals; the
biweight \citep{beaton1974fitting} and the exponential Welsch form
\citep{dennis1978techniques}, $\psi_W(r) = r \exp(-r^2 / 2c^2)$,
are the two canonical choices. We adopt Welsch for its smooth
gradient and use it in the Bayesian likelihood. It gives up a
small amount of additional efficiency for complete immunity to very
large outliers. These tools are
decades old and well understood as point estimators; their
integration into Bayesian CATE posteriors is, to our knowledge,
not standard practice.

\subsection{Heavy-tailed inference for causal effects: five paradigms}
\label{sec:related:tails}
Heavy-tailed causal inference is best understood through the
\emph{worldview} a method imposes on its tails. Five paradigms
recur in the literature, each making a different concession.

\textbf{(i) Robust M-estimation} replaces the empirical mean with
a concentration-friendly alternative --- Catoni's mean
\citep{catoni2012challenging}, median-of-means, or Huber-style
M-estimators --- while leaving the ATE/CATE estimand unchanged.
The method becomes insensitive to extreme observations at a
known efficiency cost on Gaussian data \citep{huber2009robust};
recent CATE-targeted variants (robust DR-learners,
\citet{chernozhukov2018double}-style orthogonalised scores with
robustified residuals) sit in this paradigm.

\textbf{(ii) Weight stabilisation} attacks the dual problem of
heavy tails arising from the propensity side. Overlap weights
\citep{li2018balancing}, trimming, and entropy-balancing
mitigate IPW variance explosions either by capping weights or
by redefining the target population so extreme weights vanish.
Our \texttt{use\_overlap=True} option implements the former.

\textbf{(iii) Functional replacement} changes the estimand
itself --- to a quantile, CVaR, or a full conditional
counterfactual density \citep{kallus2018policy}. The tail is
no longer a nuisance to be suppressed; it is the answer.
Distributional and quantile treatment effects are the canonical
form. Coverage of very extreme quantiles is bounded by empirical
support, where extreme value theory
\citep[EVT;][]{beirlant2006statistics} would extrapolate.

\textbf{(iv) Distribution-free intervals.} Conformal prediction
applied to ITEs
\citep{lei2021conformal, alaa2023conformal, jin2023sensitivity}
delivers finite-sample-valid prediction intervals without a tail
model. They cannot extrapolate beyond observed support, and the
recently developed sensitivity-aware variant of
\citet{jin2023sensitivity} additionally relaxes unconfoundedness.

\textbf{(v) Partial identification} concedes that the target is
not point-identified and reports sharp bounds
\citep{yadlowsky2022bounds}. Robust to arbitrarily heavy tails
within bounded-outcome assumptions, at the cost of producing a
range rather than a number.

\paragraph{Where this paper sits.}
The Bayesian X-Learner is in paradigm (i) --- it robustifies the
estimator while keeping the ATE/CATE target intact --- but adds two
moves the standard M-estimation toolkit does not. First, it returns
a \emph{full Bayesian posterior} over $\tx$ rather than a point
estimate plus sandwich SE; calibration of credible intervals under
contamination (Section~\ref{sec:experiments:coverage}) is the
empirical justification. Second, it makes the
\emph{tails-as-contamination vs tails-as-signal} distinction
operational: the same library handles paradigm-(i) contamination via
Welsch + Huber nuisance, and tail-heterogeneous $\tau$ via a
user-supplied basis $\phi(x)$ at the CATE-surface layer.

\paragraph{Connection to robust ATE pipelines.}
Industry pipelines for robust ATE estimation under heavy-tailed
revenue outcomes have used trimmed and other bounded-influence
estimators in place of plain difference-of-means or naive DR; for
example, the Trimmed Match estimator of \citet{chen2022robust}
delivers a distribution-free, robust ATE for randomised paired
geo experiments. These approaches typically terminate at a robust
point estimate (or a frequentist interval) for the ATE. Our Phase-3
Welsch posterior is complementary: it shares the bounded-influence
spirit at the regression / likelihood layer, and adds calibrated
\emph{Bayesian} UQ for the full CATE surface, supporting subgroup
contrasts, inequality probabilities, and integration over $\beta$
that a sandwich SE or bootstrap CI cannot provide.

\paragraph{Semiparametric bulk+EVT Bayesian estimators.}
An alternative approach to heavy-tailed outcomes is to model the
tail distribution directly. Semiparametric Bayesian methods based
on Dirichlet-process bulk models with a generalised Pareto (GPD)
tail component \citep{taddy2010bayesian, beirlant2006statistics}
deliver calibrated uncertainty for tail functionals (extreme
quantiles, CVaR) without requiring a contamination model. These sit
in paradigm (iii): they change the estimand to a tail-aware
functional rather than robustifying the mean. Our tail-signal basis
approach (Section~\ref{sec:experiments:tail_signal}) is simpler and
operates within paradigm (i), but a natural extension is to
integrate a GPD tail component into the Phase-3 likelihood --- we
provide preliminary results in
Section~\ref{sec:experiments:hetero_evt_phase3} and flag a
thoughtful implementation as future work.

\paragraph{RBCI and decision-theoretic calibration.}
The estimand-focused generalised-Bayes framework of
\citet{alexopoulos2025rbci} tunes the learning rate $\omega$ by minimising a
proper interval score (Winkler score) targeted at a specific causal
estimand, achieving calibration by construction for that functional.
Our trace-formula $\eta$-calibration
(Section~\ref{sec:experiments:learning_rate}) targets average
posterior-variance matching instead; a direct empirical comparison
of $\omega$-tuning vs $\eta$-trace vs $\eta^\star(a)$ on the whale
DGP is reported in Section~\ref{sec:experiments:coverage}, showing
that RBCI $\omega$-tuning trades sharpness for guaranteed
calibration while the trace formula prefers tighter intervals at
slightly sub-nominal coverage.

\paragraph{High-dimensional robust DR with CBPS/PEL.}
In high-dimensional settings, covariate-balancing propensity scores
\citep[CBPS;][]{imai2014cbps} and penalised empirical likelihood
\citep[PEL;][]{tan2020pel} provide frequentist inference under
model misspecification with bounded-influence scores. These methods
deliver root-$n$ consistent, asymptotically normal ATE estimators
without requiring correct specification of either the outcome model
or the propensity score, and their influence functions are bounded
by construction. Our method complements rather than competes with
these approaches: CBPS/PEL target scalar (or low-dimensional) ATE
functionals with frequentist guarantees, whereas our finite-
dimensional $\beta$ posterior provides a richer inferential object
(interpretability of basis coefficients, multi-contrast posterior
queries, integration over $\beta$) at the cost of a parametric
assumption on $\tau$'s functional form.

EVT-style likelihood components (paradigm-(iii)-flavoured) are not
yet fully integrated; an earlier library path
\texttt{normalize\_extremes} exposes a Hill estimator at the data
layer, but Section~\ref{sec:experiments:tail_signal} shows it is
contamination-directed rather than signal-preserving. A proper
Bayesian EVT likelihood --- generalised-Pareto tail mixture, Hill
estimator used for prior elicitation --- remains a natural
extension we flag as future work.

\paragraph{Scope and orthogonal directions.}
This paper addresses scalar continuous outcomes under cross-sectional
binary treatment. Functional outcomes \citep{salmaso2026focal} and
related longitudinal- or auxiliary-data extensions are orthogonal
directions not pursued here. Recent work on estimand-focused
generalised-Bayes tuning \citep{alexopoulos2025rbci} is closely related to our
$\eta$-calibration in
Section~\ref{sec:experiments:learning_rate}; a detailed comparison
of coverage--width trade-offs between their $\omega$-selector and
our trace-formula $\eta$ appears in
Section~\ref{sec:experiments:coverage} and
§\ref{sec:discussion}.

\paragraph{Amortised vs per-dataset Bayesian CATE.}
Recent amortised approaches to Bayesian causal inference (e.g.\
CausalFM, treating CATE as a meta-task to amortise across DGPs) and
functional-outcomes meta-learners (FOCaL,
\citealp{salmaso2026focal}) provide alternatives along orthogonal axes:
amortisation trades per-dataset MCMC for cross-task generalisation,
and FOCaL extends DR to functional outcomes. Our contribution is
per-dataset \emph{exact} robustified posterior with a principled
$\eta$-calibration step, complementary rather than competitive.
Combining tail-aware priors with amortisation, or layering Welsch
on a FOCaL-style functional-outcome regression, are natural future
directions.
